\documentclass[10pt,a4paper]{amsart}
\usepackage{amsmath,amssymb,amsthm,graphicx}
% Retained minimal preamble for the Stacks Project chapter flat.tex (More on Flatness).
% Verbatim excerpt of the chapter preamble: Xy-pic setup (native preamble lines 14-18) and the native
% theorem declarations with the mathematical macros (native preamble lines 154-217). Copied mechanically;
% nothing added, removed or redefined. The wrapper class and the amsmath/amssymb/amsthm/graphicx base
% packages are supplied by the assembler, so the class/externaldocument part of the native preamble is omitted.
% For commutative diagrams we use Xy-pic
\usepackage[all]{xy}

% We use 2cell for 2-commutative diagrams.
\xyoption{2cell}
% Theorem environments.
%
\theoremstyle{plain}
\newtheorem{theorem}[subsection]{Theorem}
\newtheorem{proposition}[subsection]{Proposition}
\newtheorem{lemma}[subsection]{Lemma}

\theoremstyle{definition}
\newtheorem{definition}[subsection]{Definition}
\newtheorem{example}[subsection]{Example}
\newtheorem{exercise}[subsection]{Exercise}
\newtheorem{situation}[subsection]{Situation}

\theoremstyle{remark}
\newtheorem{remark}[subsection]{Remark}
\newtheorem{remarks}[subsection]{Remarks}

\numberwithin{equation}{subsection}

% Macros
%
\def\lim{\mathop{\mathrm{lim}}\nolimits}
\def\colim{\mathop{\mathrm{colim}}\nolimits}
\def\Spec{\mathop{\mathrm{Spec}}}
\def\Hom{\mathop{\mathrm{Hom}}\nolimits}
\def\Ext{\mathop{\mathrm{Ext}}\nolimits}
\def\SheafHom{\mathop{\mathcal{H}\!\mathit{om}}\nolimits}
\def\SheafExt{\mathop{\mathcal{E}\!\mathit{xt}}\nolimits}
\def\Sch{\mathit{Sch}}
\def\Mor{\mathop{\mathrm{Mor}}\nolimits}
\def\Ob{\mathop{\mathrm{Ob}}\nolimits}
\def\Sh{\mathop{\mathit{Sh}}\nolimits}
\def\NL{\mathop{N\!L}\nolimits}
\def\CH{\mathop{\mathrm{CH}}\nolimits}
\def\proetale{{pro\text{-}\acute{e}tale}}
\def\etale{{\acute{e}tale}}
\def\QCoh{\mathit{QCoh}}
\def\Ker{\mathop{\mathrm{Ker}}}
\def\Im{\mathop{\mathrm{Im}}}
\def\Coker{\mathop{\mathrm{Coker}}}
\def\Coim{\mathop{\mathrm{Coim}}}

% Boxtimes
%
\DeclareMathSymbol{\boxtimes}{\mathbin}{AMSa}{"02}

%
% Macros for moduli stacks/spaces
%
\def\QCohstack{\mathcal{QC}\!\mathit{oh}}
\def\Cohstack{\mathcal{C}\!\mathit{oh}}
\def\Spacesstack{\mathcal{S}\!\mathit{paces}}
\def\Quotfunctor{\mathrm{Quot}}
\def\Hilbfunctor{\mathrm{Hilb}}
\def\Curvesstack{\mathcal{C}\!\mathit{urves}}
\def\Polarizedstack{\mathcal{P}\!\mathit{olarized}}
\def\Complexesstack{\mathcal{C}\!\mathit{omplexes}}
% \Pic is the operator that assigns to X its picard group, usage \Pic(X)
% \Picardstack_{X/B} denotes the Picard stack of X over B
% \Picardfunctor_{X/B} denotes the Picard functor of X over B
\def\Pic{\mathop{\mathrm{Pic}}\nolimits}
\def\Picardstack{\mathcal{P}\!\mathit{ic}}
\def\Picardfunctor{\mathrm{Pic}}
\def\Deformationcategory{\mathcal{D}\!\mathit{ef}}

% --- appended 2026-10-12 (Chief): the project own definitions of the reference/slogan/history
% display environments, copied mechanically from the shared Stacks preamble lines 8-12
% (\usepackage{verbatim} plus the three \newenvironment lines, which the chapter preamble keeps
% verbatim from the project preamble). Nothing added, removed or redefined: these environments
% carry no mathematical content and the project itself typesets them as comments. Without these
% three lines a counted statement whose own text uses them fails with Environment slogan undefined.
\usepackage{verbatim}
\newenvironment{reference}{\comment}{\endcomment}
\newenvironment{slogan}{\comment}{\endcomment}
\newenvironment{history}{\comment}{\endcomment}
\usepackage{cleveref}
\crefname{originallemma}{original lemma}{original lemmas}
\Crefname{originallemma}{Original lemma}{Original lemmas}
\makeatletter
\newlabel{divisors-lemma-ses-weakly-ass}{{05AN}{source}{}{sourcepointer-1}{}}
\newlabel{divisors-lemma-ses-weakly-ass@cref}{{[originallemma][0][]05AN}{[1][0][]source}{}{}{}}
\makeatother
\crefname{originallemma}{original lemma}{original lemmas}
\Crefname{originallemma}{Original lemma}{Original lemmas}
\makeatletter
\newlabel{divisors-lemma-weakly-ass-zero}{{05AP}{source}{}{sourcepointer-2}{}}
\newlabel{divisors-lemma-weakly-ass-zero@cref}{{[originallemma][0][]05AP}{[1][0][]source}{}{}{}}
\makeatother
\crefname{originallemma}{original lemma}{original lemmas}
\Crefname{originallemma}{Original lemma}{Original lemmas}
\makeatletter
\newlabel{lemma-associated-point-specializes}{{05L3}{source}{}{sourcepointer-3}{}}
\newlabel{lemma-associated-point-specializes@cref}{{[originallemma][0][]05L3}{[1][0][]source}{}{}{}}
\makeatother
\crefname{originallemma}{original lemma}{original lemmas}
\Crefname{originallemma}{Original lemma}{Original lemmas}
\makeatletter
\newlabel{lemma-bourbaki-finite-type-general-base}{{05IJ}{source}{}{sourcepointer-4}{}}
\newlabel{lemma-bourbaki-finite-type-general-base@cref}{{[originallemma][0][]05IJ}{[1][0][]source}{}{}{}}
\makeatother
\title{Injectivity of a map of finite type modules over a local base, tested on the associated points of the fibre}
\author{The Stacks Project Authors}
\date{stacks-project:flat.tex:local-snapshot-20261011; GNU FDL 1.2 or later}
\begin{document}
\maketitle
\noindent\textit{Editorial note.} All delivered mathematical text is the author's own, copied from the pinned Stacks Project chapter file flat.tex; no statement or proof step was rewritten, reordered or removed.\par\medskip
\noindent\textit{Editorial note.} Preamble provenance: the chapter's own retained minimal preamble (sources/stacks-flat/retained-preamble.tex, lines 1-85) is a verbatim excerpt of the native Xy-pic setup and of the native theorem declarations and mathematical macros, followed by the project's own reference/slogan/history display-environment definitions copied mechanically from the shared preamble; the wrapper class and the amsmath/amssymb/amsthm/graphicx base packages are supplied by the assembler, and extra\_packages is left empty because everything needed is already in the retained preamble. No macro was invented or redefined.\par\medskip
\noindent\textit{Editorial note.} Same-chapter pointers are located at their \textbackslash{}label lines inside the pinned chapter file; cross-chapter pointers are located in the project's own tag list staged in the snapshot (sources/stacks-chow/tags.txt), at the verbatim lines carrying their canonical tags.\par\medskip
\noindent\textit{Allowed standard prerequisites (editorial).}\begin{itemize}
\item Vanishing of a module with no weakly associated points, and weakly associated points of a submodule. Divisors, labels divisors-lemma-weakly-ass-zero (proof line 5566) and divisors-lemma-ses-weakly-ass (proof line 5569). Reduces injectivity to the vanishing of the weakly associated points of the kernel and places those inside the weakly associated points of the module.
\item Weakly associated points of a flat finite type module, and specialization of associated points. Same chapter flat.tex, labels lemma-bourbaki-finite-type-general-base (proof line 5571) and lemma-associated-point-specializes (proof line 5576). Shows every weakly associated point of the module is a relative associated point, and that purity makes each of those the specialization of a point of the fibre, so injectivity there propagates.
\end{itemize}\medskip
\ifdefined\hypertarget\hypertarget{sourcepointer-1}{}\fi
\noindent\textit{Source reference.} Original lemma 05AN: sources/stacks-chow/tags.txt, source lines 6856--6856. Navigation pointer for \textbackslash{}ref\{divisors-lemma-ses-weakly-ass\} used inside the counted span (line 5569); the argument that consumes it is retained in the span.\par\medskip
\ifdefined\hypertarget\hypertarget{sourcepointer-2}{}\fi
\noindent\textit{Source reference.} Original lemma 05AP: sources/stacks-chow/tags.txt, source lines 6857--6857. Navigation pointer for \textbackslash{}ref\{divisors-lemma-weakly-ass-zero\} used inside the counted span (line 5566); the argument that consumes it is retained in the span.\par\medskip
\ifdefined\hypertarget\hypertarget{sourcepointer-3}{}\fi
\noindent\textit{Source reference.} Original lemma 05L3: sources/stacks-flat/pinned-original-source.tex, source lines 5175--5175. Navigation pointer for \textbackslash{}ref\{lemma-associated-point-specializes\} used inside the counted span (line 5576); the argument that consumes it is retained in the span.\par\medskip
\ifdefined\hypertarget\hypertarget{sourcepointer-4}{}\fi
\noindent\textit{Source reference.} Original lemma 05IJ: sources/stacks-flat/pinned-original-source.tex, source lines 4079--4079. Navigation pointer for \textbackslash{}ref\{lemma-bourbaki-finite-type-general-base\} used inside the counted span (line 5571); the argument that consumes it is retained in the span.\par\medskip
\begin{lemma}
\label{lemma-injectivity-map-source-flat-pure}
Let $f : X \to S$ be a morphism of finite type.
Let $\mathcal{F}$ be a quasi-coherent sheaf of finite type on $X$.
Assume $S$ is local with closed point $s$.
Assume $\mathcal{F}$ is pure along $X_s$ and
that $\mathcal{F}$ is flat over $S$.
Let $\varphi : \mathcal{F} \to \mathcal{G}$ of quasi-coherent
$\mathcal{O}_X$-modules. Then the following are equivalent
\begin{enumerate}
\item the map on stalks $\varphi_x$ is injective for all
$x \in \text{Ass}_{X_s}(\mathcal{F}_s)$, and
\item $\varphi$ is injective.
\end{enumerate}
\end{lemma}

\begin{proof}
Let $\mathcal{K} = \Ker(\varphi)$. Our goal is to prove that
$\mathcal{K} = 0$. In order to do this it suffices to prove that
$\text{WeakAss}_X(\mathcal{K}) = \emptyset$, see
Divisors, Lemma \ref{divisors-lemma-weakly-ass-zero}.
We have
$\text{WeakAss}_X(\mathcal{K}) \subset \text{WeakAss}_X(\mathcal{F})$, see
Divisors, Lemma \ref{divisors-lemma-ses-weakly-ass}.
As $\mathcal{F}$ is flat we see from
Lemma \ref{lemma-bourbaki-finite-type-general-base}
that $\text{WeakAss}_X(\mathcal{F}) \subset \text{Ass}_{X/S}(\mathcal{F})$.
By purity any point $x'$ of $\text{Ass}_{X/S}(\mathcal{F})$
is a generalization of a point of $X_s$, and hence is the
specialization of a point $x \in \text{Ass}_{X_s}(\mathcal{F}_s)$, by
Lemma \ref{lemma-associated-point-specializes}.
Hence the injectivity of $\varphi_x$ implies the injectivity of
$\varphi_{x'}$, whence $\mathcal{K}_{x'} = 0$.
\end{proof}

\end{document}
